\documentclass{article}

% NeurIPS 2026 workshop style, double-blind submission. neurips_2026.sty bundled.
% v2 SCAFFOLD (competence-vs-performance framing). Keeps existing benchmark +
% complexity data as PRELIMINARY results; agent-sweep results are placeholders
% pending the full run. No ad-hoc "L" labels (grammatical classes are used directly).
\usepackage[dblblindworkshop]{neurips_2026}

\usepackage[utf8]{inputenc}
\usepackage[T1]{fontenc}
\usepackage{hyperref}
\usepackage{url}
\usepackage{booktabs}
\usepackage{amsmath}
\usepackage{xcolor}

% Scaffold marker for pending content.
\newcommand{\pending}[1]{\textcolor{gray}{[\textsc{pending}: #1]}}

\title{Competence is not Performance: Where Coding Agents Actually Parse\\
Along the Language-Complexity Hierarchy}

\workshoptitle{ML4Molecules: Agentic Systems for Molecular Sciences (NeurIPS 2026)}

\author{Anonymous Author(s)\\ Affiliation\\ Address\\ \texttt{email}}

\begin{document}
\maketitle

\begin{abstract}
A coding agent can emit code in a Turing-complete language, so it is routinely
\emph{assumed} able to solve information-processing tasks of arbitrary complexity.
That is a statement about expressive \emph{competence}, not realized
\emph{performance}. We make the assumption explicit and measure it, using the
first and most overlooked step of any processing chain---\textbf{parsing}---as the
probe, and real computational-chemistry output files as the test case. We place
each extraction on the language-complexity hierarchy (Chomsky, refined by Joshi's
mildly-context-sensitive band), and separate a second, orthogonal factor:
\emph{grammar-availability} (does a tested parser already exist?). Exposing tested
community parsers to the agent through an MCP harness, we vary the documentation and
scaffolding the agent receives and record its workload, chain-of-thought, and a
failure-mode taxonomy \emph{indexed by grammatical class}. \pending{headline
numbers from the full agent sweep.} Preliminary component measurements already show
the gap: on real primary-output files, structure recovery falls sharply with
grammatical class, and the deciding factor is not grammatical simplicity but whether
a grammar was \emph{supplied}. We frame the harness's role accordingly: it does not
grant the agent new expressive power (the agent could synthesize the parser itself)
but reduces the \emph{cost} of realizing it.
\end{abstract}

\section{Introduction}

\textbf{The assumption.} A coding agent emits code in a Turing-complete language, so
it is treated as able to solve arbitrarily complex processing tasks. This conflates
several claims---that the model's forward pass is Turing-complete (under idealized
precision/steps), that chain-of-thought lifts it, that it can \emph{emit} a
Turing-complete language (trivially true)---none of which establish the one that
matters: that the agent \emph{reliably produces a correct artifact} for a task of a
given complexity.

\textbf{Competence vs.\ performance.} This is Chomsky's own distinction. Expressive
\emph{competence} is an upper bound on what could be represented, not evidence it
will be found. The expressivity literature makes this explicit---capacity does not
imply generalization \citep{deletang2023chomsky}, finite precision collapses the
idealized bounds \citep{merrill2024cot}---yet agentic-systems work routinely leaves
it implicit. We treat it as a hypothesis to \emph{measure}.

\textbf{A hierarchy of complexity.} Tasks admit a principled scale: the Chomsky
hierarchy (regular $\subsetneq$ context-free $\subsetneq$ context-sensitive
$\subsetneq$ recursively enumerable), refined between context-free and
context-sensitive by Joshi's \emph{mildly context-sensitive} band and the LCFRS
fan-out ladder. We place each extraction on this scale, and keep a second factor
\emph{separate}: \emph{grammar-availability}---whether a designed, tested parser
already exists for the format (a context-free XML file is grammatically \emph{more}
complex than a regular ad-hoc log, yet far easier to extract, because its grammar is
supplied).

\textbf{Parsing as the probe.} We measure realized capability at the first,
under-studied step of any information-processing chain: parsing. Recent work grades
LLM \emph{formal-language} recognition across the Chomsky hierarchy while scaling
input length (ChomskyBench, \citealt{chomskybench}; RELIC, \citealt{relic}), our
nearest baselines---on \emph{synthetic} languages. We differ on four axes: real
heterogeneous scientific output rather than abstract strings; a
\emph{harness/documentation} manipulation; competence-vs-performance on a real
\emph{code-artifact} task (produce a working parser/pipeline, not recognize a
string); and a failure-mode taxonomy \emph{indexed by grammatical class}.

\textbf{Test case.} Real computational-chemistry / ab initio output files: hard,
heterogeneous, and where a parsing failure is a \emph{scientific} failure---the
downstream analysis silently grounds on nothing.

\section{The system under test: an MCP parser harness}

We use an open dual-mode parser harness. Community parsers are exposed both as
\textbf{MCP tools} (for interactive discovery) and as \textbf{direct Python
imports} (for production), with a \textbf{documentation/scaffolding layer}
(per-parser docs, capability metadata, prompts). It was originally built to
\emph{facilitate parser onboarding into a materials database (NOMAD)}, and later
extended toward \emph{community-wide information exchange}---parsers as shared,
discoverable, testable components. It is the right instrument here because it lets us
vary precisely the documentation and tool access the agent receives, and measure the
effect on realized performance.

\section{Method}

\paragraph{Two difficulty axes.}
(i) \emph{Source complexity} --- each format's grammatical class (regular /
context-free / mildly context-sensitive / context-sensitive), with
grammar-availability as a separate factor.
(ii) \emph{Target complexity} --- from a bounded target (fill a fixed schema) up to a
Turing-complete target (generate an executable code artifact, e.g.\ a
PySCF script), so output difficulty also spans the hierarchy.

\paragraph{The scaffolding variable (documentation arms).} The agent is given a
source file and a target schema, the parser tools, and code execution. Only the
documentation in its system prompt varies: \textsc{bare} (tools only) $\to$
\textsc{docs} (+ per-parser docs) $\to$ \textsc{full} (+ cross-parser capability map)
$\to$ \textsc{complexity} (+ the source's class and each tool's power ceiling).

\paragraph{Instrumentation.} Per run we record tool calls, turns, tokens, the agent's
\emph{chain-of-thought}, and a validated failure-mode taxonomy (wrong-tool,
silent-empty-accepted, parser-from-scratch, hallucinated-field, unit-error,
retry-loop, gave-up), indexed by grammatical class $\times$ target complexity $\times$
arm. Correctness is scored field-level against a reference built from the union of
succeeding parsers.

\section{Results}

\subsection{Preliminary: the grounding layer and the grammar-availability effect}

As a component measurement motivating the axes, we benchmarked four community parsers
on 1{,}104 real files across 16 simulation codes, scoring whether an atomic structure
was recovered. Reliability is low even on the primary-output file (Table~\ref{tab:rel}),
and failures are often \emph{silent} (an empty-but-valid object, no exception).
Re-cutting by grammatical class (Table~\ref{tab:class}) shows recovery falling as
class rises---but the effect is dominated by \emph{grammar-availability}: the one
structured (context-free) format with a supplied XML grammar is recovered every time,
while ``simpler'' regular logs with no canonical grammar are not. Grammatical class is
thus necessary but not sufficient; a supplied grammar is what makes extraction cheap.

\begin{table}[t]
\centering\small
\begin{tabular}{lcc}
\toprule
Parser & All files & Primary output \\
\midrule
A (atomistic) & 3.9\% & 23.5\% \\
B (quantum-chem) & 15.4\% & 21.6\% \\
C (regex) & 2.4\% & 7.8\% \\
D (I/O lib) & 10.5\% & 5.9\% \\
\midrule
\textbf{Oracle (any)} & \textbf{20.5\%} & \textbf{45.1\%} \\
\bottomrule
\end{tabular}
\caption{Structure-recovery reliability (preliminary component result).}
\label{tab:rel}
\end{table}

\begin{table}[t]
\centering\small
\begin{tabular}{lc}
\toprule
Source class (primary output) & Oracle recovery \\
\midrule
context-free, grammar supplied (e.g.\ XML) & 100\% \\
mildly context-sensitive (free-form log, count-agreement) & 43\% \\
context-sensitive (cross-referential / multi-file) & 0\% \\
\bottomrule
\end{tabular}
\caption{Re-cut by grammatical class: recovery falls as class rises, but the
context-free ``100\%'' reflects a \emph{supplied} grammar, not grammatical
simplicity (preliminary).}
\label{tab:class}
\end{table}

\subsection{Main: realized agent capability vs.\ complexity (overview)}

\pending{overview from the full sweep --- a single figure/table of realized
success vs.\ grammatical class $\times$ target complexity $\times$ documentation arm;
the failure-mode distribution by class; and the arm effect. Hypotheses under test:
(H1) realized success falls with class faster than a supplied-grammar baseline
predicts; (H2) the documentation arms reduce workload and wrong-tool/silent-empty
failures, concentrated in the mildly-context-sensitive band; (H3) at the
context-sensitive / Turing-complete-target cells, the residual failures are the
\emph{reasoning} modes (parser-from-scratch, compositional) that no harness removes.}
Full per-run tables and chain-of-thought transcripts appear in the SI.

\section{Discussion and outlook}

The harness's value is not expressive: the agent could synthesize the missing parser
itself (writing code is Turing-complete). Its value is to lower the \emph{cost} of
realizing that competence---supplying tested grammars, flagging emptiness, and telling
the agent when to escalate---so the offloadable (interface) failures shrink while the
irreducible (reasoning) failures remain. \pending{quantified from the sweep.} The
natural follow-up extends the target axis toward Turing-complete processing chains
(most real pipelines are Turing-complete end to end); parsing is the tractable first
step.

\section*{Limitations}

Preliminary results are a single database snapshot (52 entries, 16 codes); success is
field presence, not correctness against ground truth; parsers run with defaults; the
grammatical-class assignment is a documented ordinal partly confounded with program at
small $N$; and the main agent-sweep results are \pending{in progress}.

{\small
\bibliographystyle{plainnat}
\bibliography{refs}
}

\end{document}
